\documentclass[11pt]{article}
\usepackage[margin=25mm]{geometry}
\usepackage{fontspec}
\setmainfont{Latin Modern Roman}
\usepackage{amsmath,amssymb,mathtools}
\usepackage{unicode-math}
\setmathfont{Latin Modern Math}
\usepackage{xurl,hyperref}
\hypersetup{colorlinks=true,urlcolor=blue}
\setcounter{secnumdepth}{0}
\setlength{\parindent}{0pt}
\setlength{\parskip}{5pt}
\setlength{\emergencystretch}{3em}
\newcommand{\sref}[2]{\hyperlink{source:#1:#2}{[S#2]}}
\newcommand{\cataloguescope}[1]{\paragraph{Recorded target.}#1}
\newcommand{\R}{\mathbb R}
\newcommand{\C}{\mathbb C}
\newcommand{\N}{\mathbb N}
\newcommand{\Z}{\mathbb Z}
\newcommand{\Q}{\mathbb Q}
\newcommand{\F}{\mathbb F}
\newcommand{\eps}{\varepsilon}
\begin{document}
% UnsolvedMath ID 30000479; source catalogue code OWR-1270-002
\setcounter{section}{30000478}
\section{Quantum Zeta Functions for Mirror Calabi–Yau Threefolds}\label{30000479}
{\small UnsolvedMath ID 30000479 · Mathematical Physics}
\subsection{Definitions and mathematical statement}

\paragraph{Shared notation.}$\mathbb N=\{1,2,\ldots\}$; the standard number fields have their usual meanings. All additional conventions are stated locally.


\paragraph{Definitions and conventions.} For a smooth Calabi--Yau threefold $X$ over $\mathbb Q$ and a good prime $p$, its local zeta function encodes point counts, $Z_p(X,T)=\exp(\sum_{m\ge1}|X(\mathbb F_{p^m})|T^m/m)$; its cohomological factors are characteristic polynomials of Frobenius. An $L$-series assembles corresponding local factors over primes. A mirror pair exchanges the appropriate Hodge data. The weight motive in the report's explicit orbifold families contributes a common essential factor to the third-cohomology Frobenius polynomials of the pair.

\paragraph{Source scope and qualifications.} Quantum zeta function and quantum L-series are requested notions, not uniquely defined inputs. The report's ratio candidate has the desired duality but discards the weight-motive factor; its explicit orbifold results are context, not a universal construction for all Calabi--Yau mirror pairs.

\paragraph{Catalogue mathematical target.}
\paragraph{Catalogue wording retained for traceability.}
For mirror pairs of Calabi–Yau threefolds over $\mathbb Q$, find a satisfactory quantum zeta function and quantum $L$-series that retain the essential weight-motive factor and satisfy the expected $T$-duality.

\paragraph{Saved literature context.} The supplied catalogue's assessment dated 2026-08-21 records: The report proves restricted motive and slope-zeta statements, while later work still treats the general arithmetic and quantum zeta-function correspondence as conjectural. The assessment is retained as dated provenance; it is not a new comprehensive literature or open-status certification.
\subsection{Short English statement}
Construct arithmetic mirror invariants with the expected duality while retaining the important motive information lost by simple zeta-function ratios.
\subsection{Research attempt}
\paragraph{Model, date and investigated scope.}
OpenAI GPT-6 Astra (configured model \texttt{gpt-6-astra}), ultra reasoning; generated 11 October 2026. Proves unavoidable common-factor loss in ratio-only candidates and tests an information-preserving formal augmentation.

\paragraph{Sources and literature check.}
The complete supplied definition, including source qualifications and the dated literature assessment, was reviewed. The calculation below is an elementary restricted attempt motivated by [S2]; it does not claim novelty or a fresh literature-status determination.

\paragraph{Context and investigated scope.} The complete supplied definition, including source qualifications and the dated literature assessment, was reviewed. The calculation below is an elementary restricted attempt motivated by [S2]; it does not claim novelty or a fresh literature-status determination.

\paragraph{An information-loss calculation.} Work in the ring of formal power series over $\mathbb Q$ and take local factors $Z_X(T),Z_Y(T)$ with constant term one. If both contain a common nontrivial factor $W(T)$, say $Z_X=WA$ and $Z_Y=WB$, then their ratio is $A/B$. For every other unit $\widetilde W$, the pair $(\widetilde W A,\widetilde W B)$ has exactly the same ratio. Thus no construction depending only on that ratio can recover the common weight-motive factor.

\paragraph{A minimal algebraic augmentation.} Retain two quantities
\[
 R=Z_X/Z_Y,\qquad P=Z_XZ_Y.
\]
Exchanging the pair sends $(R,P)$ to $(R^{-1},P)$. Moreover $Z_X^2=PR$ and $Z_Y^2=P/R$, so constant-term-one square roots reconstruct the original ordered factors uniquely. To prove uniqueness and existence, the coefficient of $T^n$ in $(1+\sum_{j\geq1}a_jT^j)^2$ is $2a_n$ plus a polynomial in earlier coefficients; division by two recursively determines $a_n$. The common factor now remains as $W^2$ in $P$.

\paragraph{Why this is insufficient.} The pair $(R,P)$ is merely repackaged classical zeta data. A formal square root need not be a rational function or a motivic Euler factor; for example $\sqrt{1+T}$ has an infinite binomial series. Nor does an involution of this pair prove the source's expected physical $T$-duality or construct a global quantum $L$-series. The first gap is a geometric or physical principle selecting a new invariant with acceptable local factors while preserving the essential motive. The calculation rigorously diagnoses the ratio's loss without proposing its bookkeeping repair as the requested theory.

\paragraph{Verification and adversarial audit.}
Checked constant terms, exchange involution, square-root recursion and the distinction between formal series and rational motivic factors. Model self-audit only; no independent human proof checking. Prompts A and B were applied to this canonical record and its qualified source scope. Checked constant terms, exchange involution, square-root recursion and the distinction between formal series and rational motivic factors. Retained unresolved status and zero original-target progress; No motivic/physical principle, acceptable Euler-factor construction or quantum T-duality statement is supplied.

Coordinator checked cancellation of common formal units and the constant-term-one square-root coefficient recursion over the rational formal-series ring. This additional check was model-assisted, not independent human verification.

\paragraph{Tools, process and contributions.}
Python 3 was used for file preparation and structural checks; no numerical experiment or theorem prover was used. Contribution prompt A was followed by prompt B. GPT-6 Astra with ultra reasoning drafted the research and performed its stated self-audit; any additional coordinator contribution is identified in the verification and audit. A GPT-6-based Codex coordinator assembled the manuscript and checked repository integration; its exact variant and reasoning configuration were not exposed. No human mathematical contribution, independent human proof check or formal proof certification is supplied. Complete derivations and the consulted source references are included here; private reasoning transcripts are not published.

\paragraph{Final status and first remaining gap.}
\textbf{UNRESOLVED}. No motivic/physical principle, acceptable Euler-factor construction or quantum T-duality statement is supplied.

\section{Completion Estimate}
\noindent\textbf{COMPLETION: 0\%}
The calculation is an auxiliary or restricted attempt and establishes no newly verified portion of the full source target. No motivic/physical principle, acceptable Euler-factor construction or quantum T-duality statement is supplied.

\subsection{Sources}
\begin{itemize}
\item[\hypertarget{source:30000479:1}{S1}] ULAM.AI and the UnsolvedMath Contributors, UnsolvedMath, record 30000479 (OWR-1270-002). Catalogue curation under CC BY 4.0; definitions and English exposition added locally.
\url{https://huggingface.co/datasets/ulamai/UnsolvedMath}.
\item[\hypertarget{source:30000479:2}{S2}] Mini-Workshop: Zeta Functions, Index and Twisted K-Theory; Interactions with Physics (2006), Oberwolfach Reports, DOI 10.4171/owr/2006/21, pp. 1245--1284; official report source and mathematical context.
\url{https://doi.org/10.4171/owr/2006/21}.
\item[\hypertarget{source:30000479:3}{S3}] Official EMS Press PDF of the same Oberwolfach report; PDF pages 29--31, printed pages 1273--1275.
\url{https://ems.press/content/serial-article-files/46051}.
\item[\hypertarget{source:30000479:4}{S4}] Wan, Mirror Symmetry for Zeta Functions; reference in the saved dated catalogue literature review.
\url{https://arxiv.org/abs/math/0411464}.
\item[\hypertarget{source:30000479:5}{S5}] Zeta functions of alternate mirror Calabi-Yau families; reference in the saved dated catalogue literature review.
\url{https://arxiv.org/abs/1612.09249}.
\end{itemize}
\label{end:30000479}
\end{document}
